\chapter{H-P\(^{\mathrm{Ar}}\) Spectral Matching and Hamburger Determinacy}
\label{v4-arith:hpar-det}

\emph{Chapter~\ref{v4-arith:HPAr-vbconverse} reduced the arithmetic
Hilbert--P\'olya hypothesis H-P\(^{\mathrm{Ar}}\) to a spectral
realisation (SP) equivalent to the Connes adele-class-space
absorption-spectrum conjecture and to the Deninger regularised-determinant
axioms, and reduced the VB\(^{*}\) converse to a four-way equivalence
with Hamburger determinacy of the symmetrised shifted-zero measure.
Chapter~\ref{v4-arith:platonic-rh-completion} assembled the conditional
RH completion package and isolated the positivity/determinacy input
which would close that package. This chapter pushes on that frontier
without claiming closure. On the H-P\(^{\mathrm{Ar}}\) side, it separates
the Bost--Connes zero-mode \emph{weight-selection} problem from the
trace-matching problem: the logarithmic Berry--Keating weight is the
only survivor among a small canonical test family, and is unique under
an explicit KMS/zeta-partition normalisation, but existence of the
required trace identity remains the H-P\(^{\mathrm{Ar}}\) obstruction.
On the VB\(^{*}\) side, it audits the Carleman route: the
Bombieri--Lagarias Li-coefficient bounds and the Carleman estimate for
the shifted-zero moments are logically independent. What remains is a
precise sufficient analytic residual
\((\mathrm{R}_{\mathrm{analytic}})\): a linear upper-growth bound for
\(\mu_{2n}^{1/(2n)}\), or equivalently a positive regularised
moment measure with Carleman divergence. Both residuals are open;
the chapter produces the sharpest available reductions.}

\section{Route~3 Context}
\label{v4-arith:hpar-det:context}

\subsection{What Route~2 Settled}
\label{v4-arith:hpar-det:wave2-recap}

The route~2 structural analysis of
Chapter~\ref{v4-arith:HPAr-vbconverse} established the following
four results, each of which route~3 operates within:

\begin{enumerate}
\item \emph{Domain theorem for \(L_{0}^{\mathrm{Ar}}\).}
\Cref{v4-arith:hpvbconv:thm:L0-Friedrichs}: on the primary arithmetic
Heisenberg Fock, under non-negativity of the Bost--Connes zero-mode
weights \(\lambda(m)\geq 0\), the symmetric operator
\(L_{0}^{\mathrm{Ar},\,\mathrm{symm}}\) admits a unique Friedrichs
self-adjoint extension \(L_{0}^{\mathrm{Ar}}\), essentially self-adjoint
on the finite-span core \(\mathcal{D}_{0}\), with non-negative spectrum.
\textbf{Closed}.
\item \emph{Three-way equivalence for H-P\(^{\mathrm{Ar}}\).}
\Cref{v4-arith:hpvbconv:thm:HPAr-equivalence}: H-P\(^{\mathrm{Ar}}\)
(in the spectral-realisation form SP) is equivalent to the Connes
absorption-spectrum conjecture and to the Deninger regularised
determinant condition. \textbf{Structurally closed; inner content
open}.
\item \emph{Four-way equivalence for VB\(^{*}\) converse.}
\Cref{v4-arith:hpvbconv:thm:VBstar-converse-spectral-regularity}:
the VB\(^{*}\) converse is equivalent to Hamburger determinacy (Det) of
the symmetrised shifted-zero measure, to Hermite--Biehler operator
positivity on the de Branges space \(\mathcal{H}(E_{\xi})\) (HB+), and
to a sharp B\'aez-Duarte decay rate
\(\epsilon_{N}^{\mathrm{BD}}\asymp 1/\sqrt{\log N}\) with the
Petersson-matched constant (BD+). \textbf{Structurally closed; inner
content open}.
\item \emph{Logical independence of the two residuals.}
\Cref{v4-arith:hpvbconv:prop:HPAr-det-independent}: H-P\(^{\mathrm{Ar}}\)
(existence of a self-adjoint \(L_{0}^{\mathrm{Ar}}\) with prescribed
trace) and Hamburger determinacy (uniqueness of the spectral
measure) are logically independent. \textbf{Proved}.
\end{enumerate}

Route~3 does not repair the route~2 results; those results are correct. Route~3
performs two genuinely new operations inside that framework.

\subsection{Two Open Items Attacked Here}
\label{v4-arith:hpar-det:wave3-attack}

(A) \emph{Zero-mode weight selection for H-P\(^{\mathrm{Ar}}\).}
Route~2's \Cref{v4-arith:hpvbconv:cor:HPAr-single-residual} names the
single residual as ``the choice of Bost--Connes zero-mode weights
\(\{\lambda(m)\}_{m\in\mathbb{N}^{\times}}\) such that the combined
trace matches the zero generating function''. Route~3 tests four
canonical weight systems against the arithmetic Koszul
self-duality (P1), the arithmetic Zhu modularity (P2), and
Bost--Connes rigidity at \(\beta=1\). This is a finite stress test,
not a classification of all additive functions on
\(\mathbb{N}^{\times}\). A genuine uniqueness theorem requires the
extra KMS/zeta-partition normalisation stated below.

(B) \emph{Hamburger determinacy attack on the symmetrised shifted-zero
measure.} Route~2's
\Cref{v4-arith:hpvbconv:prop:carleman-insufficient} names the Carleman
sum as ``open, not refuted'' under a heuristic \(\mu_{2n}\asymp
C^{n}n!^{2}\). Route~3 tests the tempting Bombieri--Lagarias route and
finds the obstruction: Li-coefficient bounds do not transfer to
two-sided estimates for the positive regularised shifted-zero moments.
The usable conclusion is the classical one: Carleman's criterion would
follow from an upper bound \(\mu_{2n}^{1/(2n)}\le Cn\) for all large
\(n\), but Vol~IV does not prove that bound.

The single positivity corollary
\Cref{v4-arith:platonic:cor:single-positivity} is sufficient for RH
after the conditional Ch~37 package is supplied. Route~3 produces the
sharpest reduction the framework allows without that closure.

\subsection{Disjoint Source Catalog for Route~3}
\label{v4-arith:hpar-det:sources}

Route~2 used: Connes arXiv:math/9811068, Deninger 1994, Berry--Keating
1999, Reed--Simon, Akhiezer--Glazman, de Branges 1968/1994,
B\'aez-Duarte 2003, Burnol 2000. Route~3 draws on a complementary
catalog to preserve source disjointness:

\begin{description}
\item[Akhiezer, \emph{The Classical Moment Problem} (Oliver \& Boyd
1965).] Definitive monograph on the Hamburger, Stieltjes, Hausdorff
moment problems; \S2.4 Carleman's criterion; \S2.5 Nevanlinna
parametrisation of indeterminate solutions; \S3.1 Jacobi-matrix
formalism; \S5.6 exponential-moment estimates.
\item[Shohat--Tamarkin \emph{The Problem of Moments} (AMS 1943).] The
American companion to Akhiezer; Ch.~III.5 classical determinacy,
Ch.~III.2 Hankel determinants as sufficient tests, Ch.~VI bounded-case
theorems.
\item[Stieltjes, Ann.~Fac.~Sci.~Toulouse (1894).] The original
indeterminate moment problem; the lognormal \(\mu_{n}=e^{n^{2}/2}\)
example; baseline indeterminacy standard.
\item[Nevanlinna, Ann.~Acad.~Sci.~Fenn.~5 (1922).] Parametrisation of
the convex set of solutions to an indeterminate Hamburger moment
problem by four entire functions \(A,B,C,D\) of exponential type zero;
each indeterminate solution corresponds to a Nevanlinna function
\(\varphi\) via \((A\varphi-C)/(B\varphi-D)\).
\item[Krein, Dokl.~Akad.~Nauk SSSR 77 (1951) 93--96.] Krein's
Pontryagin-space theorem: self-adjoint operators on indefinite-metric
spaces; prototype for the indeterminate moment problem as a problem
with non-unique positive lift.
\item[Simon, Adv.~Math.~137 (1998), 82--203.] Modern treatment of
classical moment problems via Jacobi matrices: the Hamburger moment
problem with moments \(\{\mu_{n}\}\) is determinate iff the associated
Jacobi operator is essentially self-adjoint on \(c_{00}\); indeterminate
iff deficiency indices \((1,1)\).
\item[Berg--Christiansen, Math.~Ann.~322 (2002) 451--473.] Exponential
moment condition: \(\sum_{n}\mu_{2n}^{-1/(2n)}=+\infty\) implies
determinacy (Carleman); refinements via sharper generating-function
analysis.
\item[Lin, Probab.~Surv.~14 (2017) 129--190.] Modern survey of
determinacy: Carleman vs.~Krein vs.~Lin conditions; exponential-moment
comparison.
\item[Bost--Connes, Selecta Math.~1 (1995) 411--457.] Original
Bost--Connes system; \S2.2 construction of the Hecke algebra
\(C^{*}_{\mathbb{Q}}\); \S3.1 KMS states for \(\beta>1\); \S4 phase
transition at \(\beta=1\).
\item[Laca, J.~Funct.~Anal.~152 (1998) 330--378.] Laca's categorical
approach to the Bost--Connes system: characters of the
\(C^{*}(\mathbb{Q}/\mathbb{Z})\rtimes\mathbb{N}^{\times}\) crossed
product; KMS classification via induced representations.
\item[Connes--Marcolli, AMS Colloq.~Publ.~55 (2008).] Comprehensive
reference on Bost--Connes and number theory; Ch.~3 KMS states,
Ch.~4 rigidity at \(\beta=1\), Ch.~5 adele class space structure.
\item[Connes--Marcolli--Ramachandran, arXiv:math/0501424.] Explicit
description of the KMS\({}_{\beta=1}\) state as the unique symmetric
trace on the Bost--Connes system; rigidity of the zero-mode spectrum
under the Hecke action.
\item[Conrey, J.~Reine Angew.~Math.~399 (1989) 1--26.] Zero-density
theorem: at least 40\% of the non-trivial zeros of \(\zeta\) lie on the
critical line. Unconditional.
\item[Selberg, Arch.~Math.~Naturvid.~48 (1946) 89--155.] Positive
proportion of zeros on the critical line, first unconditional
zero-density result.
\item[Levinson, Adv.~Math.~13 (1974) 383--436.] 1/3 of zeros on the
line via Levinson's method.
\item[Odlyzko, 2001 (unpublished note).] Numerical verification of
\(10^{13}\) non-trivial zeros on the critical line.
\item[Bombieri--Lagarias, J.~Number Theory 77 (1999) 274--287.]
Unconditional upper bound
\(\lambda_{n}\leq n(c_{0}+\log n)/2\), \(c_{0}=2-\log(4\pi)+\gamma\), and
a matching lower bound \(\lambda_{n}\geq -(n/2)\log n-O(n)\).
\end{description}

These sources are disjoint from the route~2 catalog in the Vol~IV
sense: none of the route~2 primary derivations are cited route~3, and
none of the route~3 verifications are cited route~2.

\section{Route~3 Analysis: Six Diagnostic Cycles}
\label{v4-arith:hpar-det:ledger}

Route~3 executes six diagnostic cycles, four on weight selection
(A) and two on Hamburger determinacy (B).

\subsection{Cycle 1 (A, weight enumeration)}
\label{v4-arith:hpar-det:cycle-1}

\paragraph{Question.} Does the Berry--Keating ansatz
\(\lambda(p^{k})=k\log p\)
(\Cref{v4-arith:hpvbconv:rem:berry-keating}) constitute the unique
natural Bost--Connes zero-mode weight?

\paragraph{Analysis.} The Berry--Keating ansatz reproduces the
leading asymptotic zero density \(N(T)\sim(T/2\pi)\log(T/2\pi)\)
(Bohr--Sommerfeld quantisation of \(xp=E\)) but does not reproduce the
individual spectrum; three independent competitors exist among finite
canonical tests.

\paragraph{Sources.} Bost--Connes 1995 \S2.2 (Hecke algebra
structure); Connes--Marcolli 2008 Ch.~3 (KMS\({}_{\beta}\)
classification); Laca 1998.

\paragraph{Outcome.} Four candidate weight systems are enumerated as
\Cref{v4-arith:hpar-det:def:candidate-weights}; pairwise
P1+P2 compatibility is proved as
\Cref{v4-arith:hpar-det:thm:two-cancellations}. The KMS\({}_{\beta=1}\) state of the
Bost--Connes system is unique in the symmetric normalisation, but that
rigidity alone does not classify arbitrary additive prime-weight
systems. Trace-compatibility
\eqref{v4-arith:hpvbconv:eq:BC-trace-condition} for each candidate
is equivalent to H-P\(^{\mathrm{Ar}}\) for that candidate; no
combination closes unconditionally.

\subsection{Cycle 2 (A, uniqueness under KMS rigidity)}
\label{v4-arith:hpar-det:cycle-2}

\paragraph{Question.} Does at most one weight system satisfy
simultaneous compatibility with P1, P2, and the
H-P\(^{\mathrm{Ar}}\) trace condition?

\paragraph{Analysis.} Existence is open (it is H-P\(^{\mathrm{Ar}}\)
itself); uniqueness holds conditionally on existence, via
Connes--Marcolli rigidity.

\paragraph{Sources.} Connes--Marcolli, AMS Colloq.~Publ.~55
(2008), Thm~4.15 (uniqueness of symmetric KMS\({}_{\beta=1}\) state);
Connes--Marcolli--Ramachandran arXiv:math/0501424 (explicit form).

\paragraph{Outcome.} The KMS\({}_{\beta=1}\) state on
\(C^{*}_{\mathbb{Q}}\) is unique among symmetric states
(invariant under \(\mathrm{Gal}(\mathbb{Q}^{\mathrm{ab}}/\mathbb{Q})\))
by Connes--Marcolli 2008 Thm~4.15. The compatibility theorem
identifying the partition function and prime-generator expectations
with the canonical zeta normalisation forces the logarithmic weight
under that hypothesis; see
\Cref{v4-arith:hpar-det:thm:uniqueness-given-existence}. Existence
of the weight system, equivalent to H-P\(^{\mathrm{Ar}}\), remains open.

\subsection{Cycle 3 (A, consistency with H-P\(^{\mathrm{Ar}}\))}
\label{v4-arith:hpar-det:cycle-3}

\paragraph{Question.} Does the Berry--Keating weight
\(\lambda(p^{k})=k\log p\) reproduce the H-P\(^{\mathrm{Ar}}\) trace
identity?

\paragraph{Analysis.} The Berry--Keating partition function
\(\sum_{n\geq 1}e^{-t\log n}=\sum_{n\geq 1}n^{-t}=\zeta(t)\), multiplied
by the Dedekind-eta inverse, does not equal the zero
wave-kernel distribution;
the two sides match only after the archimedean Euler factor is
subtracted and the logarithmic-derivative expansion is inverted.

\paragraph{Sources.} Edwards, \emph{Riemann's Zeta Function}
(Academic Press 1974) \S\S1.13 and 3.3, Mellin inversion of \(\xi\)
via the logarithmic derivative.

\paragraph{Outcome.} The trace-compatibility
condition is an infinite hierarchy of constraints, not a single
scalar equation. By Mellin inversion, the hierarchy is an equality of distributions for the normal zero
operator \(\Theta_{\xi}\), equivalently
\(\det_{\infty}(s\,\mathrm{id}-\Theta_{\xi})=\xi_{\mathbb R}(s)\). The
Friedrichs Fock/BC Hamiltonian supplies the Euler partition side, not
this determinant identity. The hierarchy is identified with a regularised
Weil/Connes explicit-formula extension only after the admissibility of the
test family is supplied; see
\Cref{v4-arith:hpar-det:thm:hierarchy-condition}. No weight system is
known to close the trace identity.

\subsection{Cycle 4 (B, Hamburger determinacy reductions)}
\label{v4-arith:hpar-det:cycle-4}

\paragraph{Question.} Does Carleman's criterion
\(\sum_{n\geq 1}\mu_{2n}^{-1/(2n)}=+\infty\) hold for the
symmetrised shifted-zero moment sequence \(\{\mu_{2n}\}\)?

\paragraph{Analysis.} The moment sequence grows like
\(\mu_{2n}\asymp C^{n}n!^{2}\) under Riemann--von Mangoldt heuristics
(Chapter~\ref{v4-arith:HPAr-vbconverse}); this is exactly on the
Carleman boundary, not strictly inside.

\paragraph{Sources.} Akhiezer 1965 \S2.4 (Carleman's criterion);
Nevanlinna 1922 (parametrisation of indeterminate solutions);
Shohat--Tamarkin 1943 Ch.~III.5; Berg--Christiansen 2002
(exponential-moment refinements); Lin 2017 (modern survey).

\paragraph{Outcome.} The Riemann--von Mangoldt formula
controls zero counting, not the positivity of the regularised
moment functional. A saddle-point heuristic suggests the
Carleman-scale growth \(C^{n}n!^{2}\), but converting that heuristic into
the upper bound \(\mu_{2n}^{1/(2n)}=O(n)\) requires an additional analytic
theorem. The Hamburger determinacy problem reduces to positivity plus a
Carleman-scale upper-growth estimate on \(\mu_{2n}\); see
\Cref{v4-arith:hpar-det:thm:carleman-reformulation}. Unconditional closure
is open.

\subsection{Cycle 5 (B, attack via Li--Bombieri--Lagarias)}
\label{v4-arith:hpar-det:cycle-5}

\paragraph{Question.} Can the unconditional Bombieri--Lagarias bound
\(\lambda_{n}\leq n(c_{0}+\log n)/2\) be converted into an
unconditional estimate on \(\mu_{2n}\) sufficient for Carleman?

\paragraph{Analysis.} The Li coefficients are oscillatory sums over
zeros, whereas \(\mu_{2n}\) is a positive regularised moment. The
unconditional bound on \(\lambda_{n}\) does not
transfer to \(\mu_{2n}\) without a separate positivity/domination
theorem.

\paragraph{Sources.} Bombieri--Lagarias J.~Number Theory 77
(1999) 274--287, Theorem 2; Voros 2003 Ann.~Inst.~Fourier (refinements);
Coffey 2005 Math.~Phys.~Anal.~Geom.

\paragraph{Outcome.} Bombieri--Lagarias Thm~2 bounds the
Li sum
\(\lambda_{n}=\sum_{\rho}(1-(1-1/\rho)^{n})\) (Li 1997). This is a
highly oscillatory zero sum. It is not a bound on the absolute moments
\(\mu_{2n}\), and cancellations in the Li sum cannot be reversed into
Carleman control without an additional positivity and regularisation
theorem; see
\Cref{v4-arith:hpar-det:prop:BL-carleman-partial}. An independent
construction of a positive regularised shifted-zero moment measure with
upper-growth estimate \(\mu_{2n}^{1/(2n)}\le Cn\) remains open.

\subsection{Cycle 6 (cross-check; residual obstruction naming)}
\label{v4-arith:hpar-det:cycle-6}

\paragraph{Question.} Does the Ch~37 single positivity theorem reduce
to an analytic statement of classical form?

\paragraph{Analysis.} Comparison of the three classical reformulations
(Carleman, Nevanlinna, Akhiezer--Krein) on the symmetrised shifted-zero
measure yields a single precisely-named residual.

\paragraph{Sources.} Accumulated route~2 and route~3 sources cited above.

\paragraph{Outcome.} The route~3 reductions consolidate into a single
residual: a positive regularised interpretation of
\(\mu_{2n}=\sum_{\rho}(1/4+|\gamma_{\rho}|^{2})^{n}\) with
\(\mu_{2n}^{1/(2n)}=O(n)\). The arithmetic content required is
an analytic estimate strong enough for Carleman's criterion, together
with positivity of the regularised measure; all other ingredients
are conditional inputs from route~2 and Ch~37. See
\Cref{v4-arith:hpar-det:thm:single-analytic-residual}.

\section{H-P\(^{\mathrm{Ar}}\) Spectral Matching: Bost--Connes Zero-Mode Weight}
\label{v4-arith:hpar-det:weight-selection}

\subsection{Candidate Weights}
\label{v4-arith:hpar-det:candidate-weights}

The route~2 candidate weight system inscribed was the Berry--Keating
ansatz \(\lambda(p^{k})=k\log p\) extended multiplicatively to
\(\lambda(m)=\log m\) (\Cref{v4-arith:hpvbconv:rem:berry-keating}). We
test it against three nearby diagnostics. This is deliberately not an
exhaustive classification of all maps
\(\mathbb{N}^{\times}\to\mathbb{R}\): additive weights are determined
by arbitrary prime values \(\{\lambda(p)\}_{p}\). The point is to
separate the canonical BC/zeta generator from common false
substitutes.

\begin{definition}[four canonical zero-mode weight tests]
\label{v4-arith:hpar-det:def:candidate-weights}
\ClaimStatusDefinitional. Let \(C^{*}_{\mathbb{Q}}\) be the
Bost--Connes algebra of
\Cref{v4-arith:def:bost-connes-algebra}. Four canonical weight tests
\(\{\lambda_{\star}(m)\}_{m\in\mathbb{N}^{\times}}\) are available,
each arising from a distinct structural interpretation of the
Hecke action on the zero-mode lattice:
\begin{description}
\item[(BK)] \emph{Berry--Keating / logarithmic}:
\(\lambda_{\mathrm{BK}}(m):=\log m\). Extended multiplicatively:
\(\lambda_{\mathrm{BK}}(p^{k})=k\log p\). This is the Euler-product
generator; yields the BC KMS\({}_{\beta}\) partition function
\(\zeta(\beta)\) at \(\beta=1\).
\item[(Li-shift)] \emph{Li spectral translation}:
\(\lambda_{\mathrm{Li}}(m):=\log m\) as a semigroup weight, with the
critical-line translation \(s\mapsto s-\frac12\) applied to the
spectral parameter rather than inserted as a constant shift in the
semigroup homomorphism. A literal formula
\(\log m-\frac12\) is not additive on
\(\mathbb{N}^{\times}\) and is therefore not a BC weight.
\item[(vM)] \emph{von Mangoldt / Weil}:
\(\lambda_{\mathrm{vM}}(m):=\Lambda(m)\log m\) where
\(\Lambda(m)\) is the von Mangoldt function
(\(\Lambda(p^{k})=\log p\), zero otherwise). This is a logarithmic
derivative diagnostic, not a multiplicative semigroup weight.
\item[(FS-quad)] \emph{quadratic finite-place diagnostic}:
\(\lambda_{\mathrm{FS}}(p^{k}):=k^{2}\log p\) and
\(\lambda_{\mathrm{FS}}(m):=\sum_{p^{k_{p}}\Vert m}k_{p}^{2}\log p\).
This records the quadratic-exponent alternative suggested by
finite-place Hecke powers on perfectoid Iwahori charts
(\Cref{v4-arith:platonic:def:S2} finite-place orientation). It agrees
with (BK) on squarefree integers and differs on higher prime powers.
\end{description}
\end{definition}

\begin{remark}[scope of the four tests]
\label{v4-arith:hpar-det:rem:four-exhaust}
Under the \(\mathbb{Q}\)-positive-semigroup action of
\(\mathbb{N}^{\times}\) on the zero-mode lattice, a weight system
\(\lambda:\mathbb{N}^{\times}\to\mathbb{R}\) is compatible with the
Bost--Connes Hecke action
\(\mu_{n}\mu_{m}=\mu_{nm}\) only when
\(\lambda(nm)=\lambda(n)+\lambda(m)\). The additive case reduces to
specifying \(\lambda(p)\) for each prime \(p\), hence contains infinitely
many choices. The zeta partition normalisation
\[
 \sum_{m\ge 1} e^{-\beta\lambda(m)}=\zeta(\beta)
\]
for all \(\mathrm{Re}(\beta)>1\) forces \(e^{-\lambda(p)}=p^{-1}\)
prime-by-prime, and therefore forces (BK). Without that normalisation
there is no uniqueness theorem. The four tests above only record the
canonical logarithmic, translated, logarithmic-derivative, and
quadratic finite-place possibilities used in the attack-and-repair audit.
\end{remark}

\subsection{Pairwise Compatibility with P1 and P2}
\label{v4-arith:hpar-det:compat-P1-P2}

\begin{lemma}[P1 compatibility within the four tests]
\label{v4-arith:hpar-det:lem:P1-compat}
\ClaimStatusProvedHere. Within
\Cref{v4-arith:hpar-det:def:candidate-weights}, the additive
Berry--Keating weight is compatible with the formal
\(m\leftrightarrow m^{-1}\) functional-equation involution. The
Li-shift test is the same additive weight together with an external
critical-line translation. The von Mangoldt and quadratic finite-place
diagnostics are not multiplicative semigroup weights and therefore do
not define BC zero-mode weights satisfying P1.
\end{lemma}

\begin{proof}
For (BK),
\(\lambda_{\mathrm{BK}}(m)+\lambda_{\mathrm{BK}}(m^{-1})=0\) after
formal extension of \(\log\) to \(\mathbb{Q}_{>0}\). The Li shift is not the map
\(\log m-\frac12\) on \(\mathbb{N}^{\times}\), since such a map fails
\(\lambda(nm)=\lambda(n)+\lambda(m)\). It is the same additive weight
with the spectral parameter translated by \(1/2\). The von Mangoldt
test is supported on prime powers and the FS-quad test satisfies
\(\lambda_{\mathrm{FS}}(p^{a+b})\ne
\lambda_{\mathrm{FS}}(p^{a})+\lambda_{\mathrm{FS}}(p^{b})\) for
general \(a,b\). Hence neither is a BC semigroup weight.
\end{proof}

\begin{lemma}[P2 compatibility within the four tests]
\label{v4-arith:hpar-det:lem:P2-compat}
\ClaimStatusProvedHere. Within
\Cref{v4-arith:hpar-det:def:candidate-weights}, the Berry--Keating
weight is the only additive weight whose partition function is the
classical zeta partition function. The Li shift changes only the
external critical-line normalisation. The von Mangoldt and FS-quad
tests do not give the \(\zeta\)-partition function and hence do not
satisfy P2 in the arithmetic Zhu sense.
\end{lemma}

\begin{proof}
(BK) gives
\[
 \sum_{m\geq 1}e^{-t\log m}=\sum_{m\geq 1}m^{-t}=\zeta(t).
\]
The Li-shift convention leaves this semigroup partition function
unchanged and moves the \(1/2\) into the spectral variable. For (vM)
one obtains a logarithmic-derivative diagnostic supported on prime
powers, not a Dirichlet series with the \(\xi\)-functional equation.
For (FS-quad) one obtains
\[
 \prod_{p}\sum_{k\geq0}p^{-tk^{2}},
\]
not the Euler product for \(\zeta(t)\). These two diagnostics therefore
fail the P2 partition-function test.
\end{proof}

\begin{theorem}[four-test filter under P1+P2]
\label{v4-arith:hpar-det:thm:two-cancellations}
\ClaimStatusProvedHere. Among the four tests of
\Cref{v4-arith:hpar-det:def:candidate-weights}, simultaneous P1+P2
compatibility leaves only the Berry--Keating semigroup weight
\(\lambda(m)=\log m\), with the Li convention understood as the same
weight after external critical-line translation. This is a finite
filter, not a classification of all possible prime-weight systems.
\end{theorem}

\begin{proof}
\Cref{v4-arith:hpar-det:lem:P1-compat} and
\Cref{v4-arith:hpar-det:lem:P2-compat} eliminate the von Mangoldt and
FS-quad diagnostics as BC semigroup weights. The Li convention does not
produce a new additive weight; it is the Berry--Keating weight with the
spectral parameter centred at \(1/2\). Thus the only surviving
semigroup weight in the four-test family is (BK).
\end{proof}

\begin{remark}[what the filter proves]
\label{v4-arith:hpar-det:rem:BK-survivor}
\Cref{v4-arith:hpar-det:thm:two-cancellations} proves only the
finite-filter statement just stated. Full uniqueness requires the
zeta partition normalisation or an equivalent KMS compatibility
condition; without such a condition, arbitrary prime weights
\(\{\lambda(p)\}_{p}\) remain possible.
\end{remark}

\subsection{Uniqueness under KMS/Zeta Normalisation}
\label{v4-arith:hpar-det:uniqueness-given-existence}

The surviving canonical test is (BK). Existence with the
trace-compatibility condition remains open. Uniqueness after the partition function is fixed is unconditional.

\begin{theorem}[uniqueness under zeta partition normalisation]
\label{v4-arith:hpar-det:thm:uniqueness-given-existence}
\ClaimStatusProvedHere. Suppose
\(\lambda:\mathbb{N}^{\times}\to\mathbb{R}_{\geq0}\) is additive and
has zeta partition function
\[
 \sum_{m\geq1}e^{-\beta\lambda(m)}=\zeta(\beta)
\]
for all \(\mathrm{Re}(\beta)>1\), with Euler factors matched
prime-by-prime to the Bost--Connes Hecke generators. Then
\(\lambda(m)=\log m\). Consequently any H-P\(^{\mathrm{Ar}}\)
realisation satisfying this KMS/zeta normalisation uses the
Berry--Keating weight.
\end{theorem}

\begin{proof}
Additivity gives \(\lambda(p^{k})=k\lambda(p)\), so the partition
function factors as
\[
 \prod_{p}(1-e^{-\beta\lambda(p)})^{-1}.
\]
Prime-by-prime equality with
\(\zeta(\beta)=\prod_{p}(1-p^{-\beta})^{-1}\) gives
\(e^{-\lambda(p)}=p^{-1}\), hence \(\lambda(p)=\log p\) for every
prime. Additivity then gives \(\lambda(m)=\log m\).
\end{proof}

\begin{remark}[relation to Connes--Marcolli rigidity]
\label{v4-arith:hpar-det:rem:conditional-uniqueness}
Connes--Marcolli rigidity at \(\beta=1\) is the conceptual source of
this normalisation. Identifying an arbitrary weight with the unique
symmetric KMS state requires the prime-generator and partition-function
bridge proved in the theorem. The theorem records the exact conditional
uniqueness available in the Vol~IV path. H-P\(^{\mathrm{Ar}}\) remains
open: the Bost--Connes/Dedekind-eta trace identity is the unresolved
spectral matching problem.
\end{remark}

\subsection{The Hierarchy Condition}
\label{v4-arith:hpar-det:hierarchy}

\begin{lemma}[admissible Weil--Connes wave-kernel regularisation]
\label{v4-arith:hpar-det:lem:PW-wave-regularisation}
\ClaimStatusProvedHere. Let \(\xi=\xi_{\mathbb R}\) and write each
non-trivial zero as \(\rho=1/2+i\gamma_{\rho}\). Use the Fourier
convention
\[
h(r)=\int_{\mathbb R}\widehat h(u)e^{-iru}\,du .
\]
Let
\[
\mathsf{PW}:=\{h=\mathcal F^{-1}\psi:\psi\in C_c^\infty(\mathbb R)\}
\]
with the LF topology transported from \(C_c^\infty(\mathbb R)\). For
\(h\in\mathsf{PW}\), the Weil--Connes explicit formula gives
\begin{equation}
\label{v4-arith:hpar-det:eq:PW-weil}
\sum_{\rho}h(\gamma_{\rho})
=\int_{\mathbb R}h(r)\,dW_{\infty}(r)
-2\sum_{n\geq1}\frac{\Lambda(n)}{\sqrt n}\widehat h(\log n),
\end{equation}
in the normalisation of \eqref{v4-arith:hpar-det:eq:weil-explicit};
the prime sum is finite because \(\widehat h\) has compact support.

For \(\phi\in C_c^\infty((0,\infty))\), set
\[
h_{\phi}(r):=\int_0^\infty e^{-u/2}\phi(u)e^{-iru}\,du .
\]
Then \(h_{\phi}\in\mathsf{PW}\), and
\[
\langle\mathcal W_{\xi},\phi\rangle
:=\sum_{\rho}h_{\phi}(\gamma_{\rho})
\]
defines a distribution \(\mathcal W_{\xi}\in\mathcal D'((0,\infty))\).
Consequently the symbolic kernel \(\sum_{\rho}e^{-t\rho}\) is a
distributional wave kernel, not an ordinary heat trace.
\end{lemma}

\begin{proof}
This is the standard Paley--Wiener form of the Weil explicit formula:
compact support of \(\widehat h\) makes the prime sum finite and the
archimedean term is the distribution \(W_{\infty}\). The second
statement is obtained by substituting the compactly supported function
\(u\mapsto e^{-u/2}\phi(u)\) into the same formula. Approximation of
the point kernel uses mollifiers
\(\widehat h_{t,\epsilon}(u)=e^{-u/2}\epsilon^{-1}\eta((u-t)/\epsilon)\);
the convergence is in \(\mathcal D'((0,\infty))\), not pointwise heat
trace convergence.
\end{proof}

\begin{theorem}[trace-compatibility as regularised Weil/Connes hierarchy]
\label{v4-arith:hpar-det:thm:hierarchy-condition}
\ClaimStatusProvedHere \emph{(conditional on the H-P\(^{\mathrm{Ar}}\)
trace realisation)}. Let \(\lambda=\lambda_{\mathrm{BK}}\) be the
Berry--Keating weight of
\Cref{v4-arith:hpar-det:def:candidate-weights}. The
trace-compatibility condition
\eqref{v4-arith:hpvbconv:eq:BC-trace-condition} is the regularised
Weil/Connes trace hierarchy
\begin{equation}
\label{v4-arith:hpar-det:eq:weil-explicit}
\sum_{\rho}^{\mathrm{reg}}h_{t}(\gamma_{\rho})\;=\;
\int_{-\infty}^{\infty}h_{t}(r)\,dW(r)
\,-\,2\sum_{n\geq 1}\frac{\Lambda(n)}{\sqrt{n}}
\widehat{h}_{t}(\log n)
\end{equation}
for the wave-kernel test family associated to
\(e^{-t\Theta_{\xi}}\), after the admissible regularisation of
\Cref{v4-arith:hpar-det:lem:PW-wave-regularisation} is supplied. The
admissibility and regularisation are the trace-realisation
bridge required for H-P\(^{\mathrm{Ar}}\); the classical Weil formula
applies only to admissible test functions.
\end{theorem}

\begin{proof}
Under the BK weight system, the Friedrichs operator acts on
\(\delta_{m}\) by eigenvalue \(\log m\). The regularised trace of
\(e^{-tL_{0}^{\mathrm{Ar}}}\) on the zero-mode lattice is
\(\sum_{m\geq 1}e^{-t\log m}=\sum_{m\geq 1}m^{-t}=\zeta(t)\), convergent
for \(t>1\) and meromorphically continued.

Combining with the archimedean Dedekind-eta inverse product
\(\prod_{i\geq 1}(1-e^{-ti})^{-1}\) and subtracting the archimedean
descendant contribution (Tate-pole and archimedean terms) via the
Bost--Connes regularisation
(\Cref{v4-arith:hpvbconv:prop:BC-trace-compatibility}), one obtains
\begin{equation*}
\mathrm{tr}_{\mathrm{BC}}(e^{-tL_{0}^{\mathrm{Ar}}})
=\zeta(t)\cdot\eta(it)^{-1}\cdot(\text{regularisation}),
\end{equation*}
with the archimedean terms removed by the Bost--Connes/Dedekind-eta
prescription. The Hadamard product for \(\xi_{\mathbb R}\) rewrites
the zero side as a regularised zero trace. The classical Weil explicit
formula supplies this identity for the admissible Paley--Wiener test
spaces of \Cref{v4-arith:hpar-det:lem:PW-wave-regularisation}. The
Vol~IV condition is the additional assertion that the
Bost--Connes/Fock trace is carried to that distributional extension.
Therefore the theorem is a reformulation of the trace-compatibility
hierarchy under H-P\(^{\mathrm{Ar}}\), not a proof of
H-P\(^{\mathrm{Ar}}\).
\end{proof}

\begin{corollary}[the hierarchy is the Connes absorption-spectrum bridge]
\label{v4-arith:hpar-det:cor:hierarchy-is-weil}
\ClaimStatusProvedHere \emph{(conditional on the H-P\(^{\mathrm{Ar}}\)
trace realisation)}. Under the BK weight selection, the
H-P\(^{\mathrm{Ar}}\) trace-compatibility condition is the
Connes--Meyer--Weil absorption-spectrum bridge for the
Bost--Connes/Weil wave-distribution family. The classical Weil explicit
formula applies to admissible Paley--Wiener test functions; the
admissibility and trace-identification of the BC heat-trace family is
the additional content.
\end{corollary}

\begin{proof}
The previous theorem identifies the needed equality with the
regularised explicit-formula extension for the BC/Weil wave family.
By \Cref{v4-arith:hpvbconv:cor:HPAr-connes-equivalence}, this is the
same structural condition as the Connes absorption-spectrum
realisation. The admissibility and trace-identification step is
precisely what H-P\(^{\mathrm{Ar}}\) asserts.
\end{proof}

\begin{remark}[weight selection does not close H-P\(^{\mathrm{Ar}}\)]
\label{v4-arith:hpar-det:rem:weight-does-not-close}
\Cref{v4-arith:hpar-det:thm:uniqueness-given-existence} and
\Cref{v4-arith:hpar-det:cor:hierarchy-is-weil} together state: if
H-P\(^{\mathrm{Ar}}\) holds, the weight is uniquely Berry--Keating
under the zeta/KMS normalisation. Whether the Berry--Keating
weight satisfies the trace compatibility is the
regularised Connes absorption-spectrum problem. The reduction is
tight: route~3 determines the canonical weight under explicit
normalisation, but existence itself remains a classical open problem.
\end{remark}

\section{Hamburger Determinacy Attack}
\label{v4-arith:hpar-det:determinacy}

\subsection{The Symmetrised Shifted-Zero Measure}
\label{v4-arith:hpar-det:symmetrised-measure}

\begin{definition}[symmetrised shifted-zero moment sequence]
\label{v4-arith:hpar-det:def:shifted-moments}
\ClaimStatusDefinitional. Let \(\rho=1/2+i\gamma_{\rho}\) denote a
non-trivial zero of \(\xi\), written with \(\gamma_{\rho}\in\mathbb{C}\)
(real for every zero iff RH). The \emph{symmetrised shifted-zero moment
sequence} is
\begin{equation}
\label{v4-arith:hpar-det:eq:mu-2n}
\mu_{2n}\;:=\;\sum_{\rho}\Bigl(\tfrac{1}{4}+\gamma_{\rho}^{2}\Bigr)^{n}
\end{equation}
Bost--Connes regularised. Under RH, \(\gamma_{\rho}\in\mathbb{R}\) and
the unregularised terms are real and non-negative. Outside RH the
expression is not automatically a positive moment sequence; positivity
of the regularised functional is part of
\((\mathrm{R}_{\mathrm{analytic}})\).
\end{definition}

\begin{remark}[relation to the symmetrised zero measure of route~2]
\label{v4-arith:hpar-det:rem:measure-relation}
The symmetrised zero measure \(\nu\) of
\Cref{v4-arith:vb:def:zero-measure} is the pushforward of the zero
set under the symmetrisation map
\(\rho\mapsto \tfrac{1}{4}+\mathrm{Im}(\rho)^{2}\).
On the critical line \(\mathrm{Re}(\rho)=\tfrac{1}{2}\), this
equals \(|\rho|^{2}\). The shifted moment \(\mu_{2n}\) of
\Cref{v4-arith:hpar-det:def:shifted-moments} is the \(n\)-th moment of
this measure; odd moments vanish by symmetry. Determinacy of
\(\nu\) is equivalent to determinacy of \(\{\mu_{2n}\}_{n\geq 0}\)
under the standard Stieltjes--Hamburger dictionary.
\end{remark}

\subsection{Independence of Li-Coefficient Bounds and Carleman Moments}
\label{v4-arith:hpar-det:carleman-BL}

\begin{proposition}[Li-coefficient bounds and Carleman moments are independent]
\label{v4-arith:hpar-det:prop:BL-carleman-partial}
\ClaimStatusProvedHere. The Bombieri--Lagarias bounds for Li
coefficients and the Carleman estimate for the
regularised moments \(\mu_{2n}\) are logically independent.
The Bombieri--Lagarias upper bound \(|\lambda_n| = O(n\log n)\)
leaves open the implication
\[
 |\lambda_{n}|=O(n\log n)
 \quad\Longrightarrow\quad
 \mu_{2n}^{1/(2n)}=O(n):
\]
a separate positivity and domination theorem is required.
\end{proposition}

\begin{proof}
The Li coefficient is the oscillatory zero sum
\[
 \lambda_{n}=\sum_{\rho}\left(1-\left(1-\frac1{\rho}\right)^{n}\right).
\]
The shifted-zero moment relevant to Carleman is instead the positive
regularised quantity
\[
 \mu_{2n}=\sum_{\rho}^{\mathrm{reg}}\left(\frac14+\gamma_{\rho}^{2}\right)^{n}.
\]
An estimate for the oscillatory Li sum gives no bound on the absolute
or positive regularised \(2n\)-th moment without a separate positivity
and domination theorem. Reversing the cancellations in the Li sum
would itself be a new spectral-positivity input. Hence the
Bombieri--Lagarias theorem is evidence for the Li/RH reformulation but
not a Carleman proof.
\end{proof}

\begin{remark}[direction of the Carleman bound]
\label{v4-arith:hpar-det:rem:carleman-not-decided}
A lower bound on
\(\mu_{2n}\) would give an upper bound on
\(\mu_{2n}^{-1/(2n)}\), not the lower bound on
\(\mu_{2n}^{-1/(2n)}\) that Carleman's divergence criterion requires.
Carleman's criterion is driven by upper growth of the
moments, together with positivity of the representing measure.
\end{remark}

\subsection{Tightest Conditional Reduction}
\label{v4-arith:hpar-det:tightest-reduction}

\begin{theorem}[Hamburger determinacy from linear upper growth]
\label{v4-arith:hpar-det:thm:carleman-reformulation}
\ClaimStatusProvedHere. Suppose the symmetrised shifted-zero moments
come from a positive regularised Hamburger moment functional and
there is a constant \(C>0\) such that
\begin{equation}
\label{v4-arith:hpar-det:eq:sharp-asymptotic}
 \mu_{2n}^{1/(2n)}\le Cn
\end{equation}
for all sufficiently large \(n\). Then Carleman's criterion holds:
\[
 \sum_{n\ge1}\mu_{2n}^{-1/(2n)}=+\infty.
\]
Consequently the associated Hamburger moment problem is determinate.
\end{theorem}

\begin{proof}
The upper-growth hypothesis gives
\[
 \mu_{2n}^{-1/(2n)}\ge \frac{1}{Cn}
\]
for all large \(n\). The harmonic series diverges, so Carleman's
criterion applies. By the classical Carleman theorem for Hamburger
moments, the positive moment functional is determinate.
\end{proof}

\begin{corollary}[unconditional inputs and the Carleman gap]
\label{v4-arith:hpar-det:cor:unconditional-gap}
\ClaimStatusProvedHere. The current unconditional inputs of this
volume leave open the hypotheses of
\Cref{v4-arith:hpar-det:thm:carleman-reformulation}: the required positive regularised moment functional and the
linear upper-growth estimate \(\mu_{2n}^{1/(2n)}=O(n)\) are both open.
\end{corollary}

\begin{proof}
\Cref{v4-arith:hpar-det:prop:BL-carleman-partial} blocks the only
candidate unconditional shortcut in the present chapter. The
Riemann--von Mangoldt counting formula controls the distribution of
zeros by height, but by itself it does not define the required
positive regularised moment functional nor bound the regularised
moments after off-line and archimedean terms are included. These are
exactly the missing inputs.
\end{proof}

\begin{theorem}[the single analytic residual]
\label{v4-arith:hpar-det:thm:single-analytic-residual}
\ClaimStatusProvedHere. Under the route~2 structural reductions and
the route~3 analytic reductions, the Hamburger-determinacy component
of the Vol~IV RH path reduces to the following sufficient analytic
residual:
\begin{equation}
\label{v4-arith:hpar-det:eq:single-residual}
(\mathrm{R}_{\mathrm{analytic}})\quad
\text{\(\{\mu_{2n}\}\) is a positive regularised Hamburger moment
sequence and }\limsup_{n\to\infty}\frac{\mu_{2n}^{1/(2n)}}{n}<\infty.
\end{equation}
This statement implies Carleman's criterion, hence Hamburger
determinacy, hence the VB\(^{*}\) converse in the structural framework
of Chapter~\ref{v4-arith:HPAr-vbconverse}.
\end{theorem}

\begin{proof}
The implication
\((\mathrm{R}_{\mathrm{analytic}})\Rightarrow\) Carleman is
\Cref{v4-arith:hpar-det:thm:carleman-reformulation}. The
Carleman-to-determinacy implication is the classical theorem of
Akhiezer/Shohat--Tamarkin. The structural equivalence between
determinacy and the VB\(^{*}\) converse is
\Cref{v4-arith:hpvbconv:thm:VBstar-converse-spectral-regularity}.
\end{proof}

\begin{remark}[scope of \(\mathrm{R}_{\mathrm{analytic}}\)]
\label{v4-arith:hpar-det:rem:residual-content}
\((\mathrm{R}_{\mathrm{analytic}})\) is a one-sided upper-growth
condition, distinct from a two-sided asymptotic. Its relation to RH is
open: the condition may be strictly weaker, or RH-equivalent; the
chapter takes no position. It supplies the exact
positivity and upper-growth information that Carleman's
criterion requires. Standard zero-density theorems count zeros in regions;
the residual controls a regularised moment functional.
\end{remark}

\begin{remark}[status of the determinacy residual]
\label{v4-arith:hpar-det:rem:residual-vs-RH}
\((\mathrm{R}_{\mathrm{analytic}})\) is a sufficient condition for the Vol~IV
determinacy step, open as an unconditional statement. It identifies the exact analytic shape of
the positivity theorem required; neither the classical Weil formula
nor the Bombieri--Lagarias bounds supply it.
\end{remark}

\section{Consequence for Ch~37 RH Completion}
\label{v4-arith:hpar-det:ch37-consequence}

\begin{theorem}[route~3 consequence for the Ch~37 single positivity
corollary]
\label{v4-arith:hpar-det:thm:ch37-consequence}
\ClaimStatusProvedHere. The Ch~37 single positivity corollary
\Cref{v4-arith:platonic:cor:single-positivity} reduces, under
route~3 analysis, to the conjunction:
\begin{enumerate}
\item the Berry--Keating weight is forced by the additive
KMS/zeta-partition normalisation
(\Cref{v4-arith:hpar-det:thm:uniqueness-given-existence}).
\item H-P\(^{\mathrm{Ar}}\) holds with the Berry--Keating weight,
meaning that the Bost--Connes trace is identified with the
regularised Weil/Connes trace hierarchy
(\Cref{v4-arith:hpar-det:thm:hierarchy-condition}).
\item \((\mathrm{R}_{\mathrm{analytic}})\): the shifted-zero moments
form a positive regularised Hamburger moment sequence with
\(\mu_{2n}^{1/(2n)}=O(n)\)
(\Cref{v4-arith:hpar-det:thm:single-analytic-residual}).
\end{enumerate}
Under the conjunction of these three, plus the conditional P1/P3
inputs named in Chapter~\ref{v4-arith:platonic-rh-completion}, the
Ch~37 corollary closes.
\end{theorem}

\begin{proof}
(1) provides the uniquely-determined candidate operator
\(L_{0}^{\mathrm{Ar}}=L_{0}^{\mathrm{Ar}}(\lambda_{\mathrm{BK}})\).
(2) is precisely the trace-compatibility condition which, via
\Cref{v4-arith:hpvbconv:thm:HPAr-equivalence}, asserts
H-P\(^{\mathrm{Ar}}\). (3), by
\Cref{v4-arith:hpar-det:thm:single-analytic-residual} asserts
Hamburger determinacy of the symmetrised shifted-zero measure;
by
\Cref{v4-arith:hpvbconv:thm:VBstar-converse-spectral-regularity},
this supplies the VB\(^{*}\) converse. Combined with the domain theorem,
P2, and the conditional P1/P3 inputs recorded in Ch~37,
the Ch~37 single positivity corollary and its consequence RH hold.
\end{proof}

\begin{corollary}[structure of the route~3 reduction]
\label{v4-arith:hpar-det:cor:wave3-structure}
\ClaimStatusProvedHere. Route~3 reduces the Ch~37 single positivity
theorem to three named residuals:
\begin{description}
\item[(\(\mathrm{W}\))] \emph{Weight normalisation}: an additive
Bost--Connes weight with zeta partition function and prime-by-prime
KMS compatibility. Under (W), \(\lambda=\lambda_{\mathrm{BK}}\).
\item[(\(\mathrm{H}\))] \emph{H-P trace bridge}: the
Bost--Connes/Dedekind-eta heat trace is the admissible regularised
Weil/Connes zero trace. Under (H), H-P\(^{\mathrm{Ar}}\) holds for the
BK operator.
\item[(\(\mathrm{R}_{\mathrm{analytic}}\))] \emph{Carleman positivity}:
the shifted-zero moments form a positive regularised Hamburger moment
sequence and satisfy \(\mu_{2n}^{1/(2n)}=O(n)\). Under
\((\mathrm{R}_{\mathrm{analytic}})\), Hamburger determinacy holds.
\end{description}
All three together close the residual H-P/VB part of the Ch~37 path;
they still sit on top of the conditional P1/P3 inputs stated there.
\end{corollary}

\begin{proof}
Direct from \Cref{v4-arith:hpar-det:thm:ch37-consequence}. (W) fixes
the candidate operator; (H) supplies spectral matching; and
\((\mathrm{R}_{\mathrm{analytic}})\) supplies determinacy. Omitting any
one leaves either no fixed operator, no zero trace, or no VB\(^{*}\)
converse.
\end{proof}

\section{Residual Obstructions}
\label{v4-arith:hpar-det:residuals}

After the route~3 reductions, the precise remaining gaps in the
Ch~37 RH completion path are:

\begin{enumerate}
\item \textbf{(W) Weight normalisation}. The BK weight is forced once
the additive zeta partition function is fixed prime-by-prime. The
remaining task is to prove that the Vol~IV Friedrichs construction
really lands in this KMS/zeta normalisation.
\item \textbf{(H) H-P trace bridge}. The open theorem is the
admissible regularised extension identifying the BC/Dedekind-eta heat
trace with the Connes zero trace; the classical Weil explicit formula
applies to admissible test functions and is a step toward, not a
substitute for, this identification.
\item \textbf{\((\mathrm{R}_{\mathrm{analytic}})\) Carleman positivity}.
The required analytic input is a positive regularised shifted-zero
moment functional with \(\mu_{2n}^{1/(2n)}=O(n)\); this estimate
lies beyond the reach of the Bombieri--Lagarias Li-coefficient bounds
by \Cref{v4-arith:hpar-det:prop:BL-carleman-partial}.
\item \textbf{Integration with the Ch~37 package}. The RH completion
package of \Cref{v4-arith:platonic:def:RH-package} requires P1, P2,
P3, H-P\(^{\mathrm{Ar}}\), and Hamburger determinacy together. P1 is
conditional on arithmetic Positselski (scope 30--50 pages); P2
unconditional; P3 conditional on Koszul--Petersson compatibility
(scope 10--20 pages); H-P\(^{\mathrm{Ar}}\) open
(RH-equivalent); determinacy reduces to
\((\mathrm{R}_{\mathrm{analytic}})\).
\end{enumerate}

\begin{remark}[route~3 contributions]
\label{v4-arith:hpar-det:rem:wave3-contribution}
H-P\(^{\mathrm{Ar}}\) spectral matching and Hamburger determinacy
both remain RH-equivalent or RH-adjacent open problems.
Route~3 establishes:
\begin{enumerate}
\item \emph{Conditional uniqueness of the weight system}:
the Berry--Keating \(\lambda(m)=\log m\) is forced by additive
zeta/KMS normalisation.
\item \emph{Explicit identification of the trace-compatibility
hierarchy}: the infinite sequence of conditions is the regularised
Weil/Connes trace bridge; the classical Weil formula applies
to admissible test functions and is a step in the bridge.
\item \emph{Independence of Li-coefficient bounds and Carleman
moments}: the Bombieri--Lagarias bounds leave the Carleman estimate
open by \Cref{v4-arith:hpar-det:prop:BL-carleman-partial}.
\item \emph{The single analytic residual
\((\mathrm{R}_{\mathrm{analytic}})\)}: Hamburger determinacy reduces
to positivity plus linear upper growth for the regularised moments.
\item \emph{Consistency with Ch~37}: the three residuals (W), (H),
\((\mathrm{R}_{\mathrm{analytic}})\) jointly close the H-P/VB residual
after the conditional P1/P3 inputs are supplied.
\end{enumerate}
\end{remark}

\section{Summary}
\label{v4-arith:hpar-det:summary}

\begin{enumerate}
\item \textbf{Weight selection is uniquely determined after
normalisation}
(\Cref{v4-arith:hpar-det:thm:uniqueness-given-existence}). Among four
canonical tests, only the Berry--Keating \(\lambda(m)=\log m\)
survives as a BC semigroup weight with zeta partition function
(\Cref{v4-arith:hpar-det:thm:two-cancellations}).
\item \textbf{The trace-compatibility hierarchy is the regularised
Weil/Connes bridge}
(\Cref{v4-arith:hpar-det:thm:hierarchy-condition},
\Cref{v4-arith:hpar-det:cor:hierarchy-is-weil}). The
H-P\(^{\mathrm{Ar}}\) trace identity needs an admissible extension of
the explicit formula to the BC heat-trace family.
\item \textbf{Hamburger determinacy reduces to a single analytic
positivity residual}
(\Cref{v4-arith:hpar-det:thm:single-analytic-residual}). Under
\((\mathrm{R}_{\mathrm{analytic}})\), the moments are positive and
satisfy \(\mu_{2n}^{1/(2n)}=O(n)\), so Carleman's criterion holds.
\item \textbf{Li-coefficient bounds and Carleman moments are
independent}
(\Cref{v4-arith:hpar-det:prop:BL-carleman-partial}).
Li-coefficient bounds are oscillatory zero-sum bounds; a separate
positivity and domination theorem is required for Carleman.
\item \textbf{Three precisely-named residuals}
(\Cref{v4-arith:hpar-det:cor:wave3-structure}): weight
normalisation (W), H-P trace bridge (H), and Carleman positivity
\((\mathrm{R}_{\mathrm{analytic}})\).
\item \textbf{Scope of the reduction}. H-P\(^{\mathrm{Ar}}\) reduces to
the Connes absorption-spectrum conjecture (RH-equivalent); Hamburger
determinacy reduces to \((\mathrm{R}_{\mathrm{analytic}})\)
(RH-adjacent); no new RH-equivalent positivity language beyond the
route~2 catalog of Det, HB+, BD+ is produced.
\item \textbf{The sharpest available reduction}. The three residuals
(W), (H), \((\mathrm{R}_{\mathrm{analytic}})\) are each precisely
stated; their conjunction, together with the conditional P1/P3 inputs,
closes the conditional Vol~IV RH path.
\end{enumerate}

The open problems are: prove the
Bost--Connes absorption-spectrum bridge and prove positivity plus
Carleman-scale upper growth for the regularised shifted-zero moments.
